% Irodov 用户PDF页97--98，原书印刷页88--89；第2.7节220--229题。
% 使用硬球碰撞平均自由程 lambda=(sqrt2 pi d^2 n)^-1。
\thermotransportproblem{220}{88}
求气体分子的自由程长度：
(a) 超过平均自由程 $\lambda$ 的分数；
(b) 位于 $\lambda$ 与 $2\lambda$ 之间的分数。
\exsolution
硬球模型下分子从最近一次碰撞后，
在路程 $s$ 内未碰撞的存活概率为
$P(L>s)=e^{-s/\lambda}$。
(a) 直接代入 $s=\lambda$：
\[
\boxed{P(L>\lambda)=e^{-1}\simeq0.368}.
\]
(b) 区间概率等于两次存活概率之差：
\[
\boxed{P(\lambda\le L\le2\lambda)
=e^{-1}-e^{-2}\simeq0.233}.
\]
\excheck 自由程的平均值是 $\lambda$，
但其分布不是所有分子都飞过恰好 $\lambda$，
而是具有指数长尾。

\thermotransportproblem{221}{88}
一束狭窄分子束射入充满低压气体的容器。
沿传播方向经过距离 $\Delta\ell$ 后，
分子束强度降低为原来的 $1/\eta$。
求束流分子的平均自由程 $\lambda$。
\exsolution
未发生碰撞的分子数按
$I(\ell)=I_0e^{-\ell/\lambda}$ 衰减。
由
$I(\Delta\ell)/I_0=1/\eta$：
\[
e^{-\Delta\ell/\lambda}=\eta^{-1}.
\]
两边取自然对数：
$-\Delta\ell/\lambda=-\ln\eta$。
因此
\[
\boxed{\lambda=\frac{\Delta\ell}{\ln\eta}},
\quad\eta>1.
\]
\excheck 这里测到的是\emph{未碰撞的直束强度}，
要求散射分子不再被计入狭束探测器。

\thermotransportproblem{222}{88}
若气体分子在短时间
$\mathrm dt$ 内发生碰撞的概率为
$\alpha\,\mathrm dt$，其中碰撞率
$\alpha$ 为常数，求：
(a) 在时间 $t$ 内完全没有碰撞的概率；
(b) 两次相继碰撞间的平均等待时间。
\exsolution
设无碰撞存活概率为 $P_0(t)$。
在下一微时段，
$P_0(t+\mathrm dt)=P_0(t)(1-\alpha\mathrm dt)$，
故
$\mathrm dP_0/\mathrm dt=-\alpha P_0$。
用初态 $P_0(0)=1$，
得
\[
\boxed{P_0(t)=e^{-\alpha t}}.
\]
等待时间的概率密度是
$f(t)=-P_0'(t)=\alpha e^{-\alpha t}$，
平均值：
\[
\boxed{\bar\tau=\int_0^\infty
t\alpha e^{-\alpha t}\,\mathrm dt=\frac1\alpha}.
\]
\excheck $\alpha$ 的单位是 $\mathrm{s^{-1}}$，
而 $\bar\tau$ 的单位是秒。

\thermotransportproblem{223}{88}
求氮气分子的平均自由程 $\lambda$
及平均碰撞间隔 $\tau$：
(a) 标准状况；
(b) 温度 $0^\circ\mathrm C$，
压强 $p=1.0\,\mathrm{nPa}$ 的极高真空。
\exsolution
硬球模型中
\[
\lambda=\frac1{\sqrt2\,\pi d^2n}
=\frac{k_BT}{\sqrt2\,\pi d^2p},
\quad
\bar v=\sqrt{\frac{8RT}{\pi M}},
\quad \tau=\frac\lambda{\bar v}.
\]
取氮气分子有效直径
$d\simeq0.37\,\mathrm{nm}$、
$M\simeq0.028\,\mathrm{kg/mol}$、
$T\simeq273\,\mathrm K$。
(a) 标准压强
$p_0\simeq1.013\times10^5\,\mathrm{Pa}$，
得到数量级
\[
\boxed{\lambda_0\simeq0.06\,\mathrm{\mu m}},
\qquad
\boxed{\tau_0\simeq0.13\,\mathrm{ns}}.
\]
(b) 温度不变，自由程与压强成反比：
\[
\frac{\lambda_b}{\lambda_0}=\frac{p_0}{p_b}
\simeq1.0\times10^{14}.
\]
因此
\[
\boxed{\lambda_b\simeq6\times10^6\,\mathrm m
=6\,\mathrm{Mm}},\qquad
\boxed{\tau_b\simeq1.4\times10^4\,\mathrm s
\simeq3.8\,\mathrm h}.
\]
\excheck 本版答案经 PDF 文本提取后显示“6 mm”，
但与 $\tau=3.8\,\mathrm h$ 及
自由程反比压强的关系不相容；
应理解为 $6\,\mathrm{Mm}$（百万米）量级。
稀薄气体中该自由程远超通常容器尺寸，
真实分子大多先撞壁。

\thermotransportproblem{224}{88}
标准状况下，
氮气分子的平均自由程是
气体分子之间典型平均距离的多少倍？
\exsolution
若数密度为 $n$，
典型分子间距定义为
$a=n^{-1/3}$。
硬球分子自由程
$\lambda=1/(\sqrt2\pi d^2n)$，
因此比值
\[
\boxed{\frac{\lambda}{a}
=\frac{n^{-2/3}}{\sqrt2\pi d^2}}.
\]
标准状况
$n=p_0/(k_BT_0)
\simeq2.69\times10^{25}\,\mathrm{m^{-3}}$。
取有效氮分子直径
$d\simeq0.37\,\mathrm{nm}$，
有 $a\simeq3.34\,\mathrm{nm}$、
$\lambda\simeq60\,\mathrm{nm}$，
于是
\[
\boxed{\lambda/a\simeq18}.
\]
\excheck 平均自由程与分子间距不是同一个物理量；
气体虽有相邻分子，
碰撞还取决于分子的硬球碰撞截面。

\thermotransportproblem{225}{88}
一气体在标准状况下，
范德瓦耳斯排除体积参数为
$b=40\,\mathrm{mL/mol}$。
求气体分子的平均自由程。
\exsolution
硬球分子的范德瓦耳斯参数与有效直径 $d$
满足
\[
b=4N_A\left(\frac{\pi d^3}{6}\right)
=\frac{2\pi}3N_A d^3.
\]
因此
\[
d=\left(\frac{3b}{2\pi N_A}\right)^{1/3}.
\]
由硬球平均自由程
$\lambda=k_BT_0/(\sqrt2\pi d^2p_0)$，
代入 $d$：
\[
\boxed{\lambda=
\frac{k_BT_0}{\sqrt2\pi p_0}
\left(\frac{2\pi N_A}{3b}\right)^{2/3}}.
\]
用
$b=40\,\mathrm{mL/mol}=4.0\times10^{-5}\,\mathrm{m^3/mol}$，
$T_0\simeq273.15\,\mathrm K$，
\[
\boxed{\lambda\simeq8.4\times10^{-8}\,\mathrm m
=84\,\mathrm{nm}}.
\]
\excheck 需注意 $b$ 是\emph{每摩尔}的排除体积，
不是单个分子的几何体积。

\thermotransportproblem{226}{88}
声波在标准状况的氮气中传播。
求其波长与氮分子平均自由程相等时的频率。
\exsolution
声速在理想气体中为
\[
c_s=\sqrt{\frac{\gamma RT}{M}},
\quad \gamma\simeq\frac75.
\]
波长等于自由程 $\lambda$ 时
$f=c_s/\lambda$。
由 $\lambda=k_BT/(\sqrt2\pi d^2p)$
和 $R=N_Ak_B$：
\[
\boxed{f=
\frac{\sqrt2\pi d^2p_0}{k_BT_0}
\sqrt{\frac{\gamma RT_0}{M}}
=\frac{\pi d^2p_0N_A}{RT_0}
\sqrt{\frac{2\gamma RT_0}{M}}}.
\]
取 $d\simeq0.37\,\mathrm{nm}$、
$T_0\simeq273\,\mathrm K$：
\[
\boxed{f\simeq5.5\,\mathrm{GHz}}.
\]
\excheck 当波长与分子自由程同量级时，
连续介质的通常声学近似接近失效；
数值是把宏观声速外推到该尺度得到的形式估计。

\thermotransportproblem{227}{88}
温度 $0^\circ\mathrm C$ 的氧气
密封在特征尺度
$\ell=10\,\mathrm{mm}$ 的容器中。
求：(a) 分子平均自由程
$\lambda>\ell$ 所需压强范围；
(b) 临界压强处氧气分子数密度与平均间距。
\exsolution
硬球平均自由程
$\lambda=k_BT/(\sqrt2\pi d^2p)$。
令 $\lambda=\ell$，
得到临界压强
\[
\boxed{p_c=\frac{k_BT}{\sqrt2\pi d^2\ell}},
\quad
\lambda>\ell\iff p<p_c.
\]
取氧分子直径
$d\simeq0.35\,\mathrm{nm}$、
$T=273\,\mathrm K$，
\[
\boxed{p_c\simeq0.7\,\mathrm{Pa}}.
\]
临界状态数密度
$n_c=p_c/(k_BT)$，
平均间距
$a=n_c^{-1/3}$，
\[
\boxed{n_c\simeq2\times10^{14}\,
\mathrm{cm^{-3}}},\qquad
\boxed{a\simeq0.2\,\mathrm{\mu m}}.
\]
\excheck 条件为自由程超过\emph{特征几何尺度}，
不是要让分子间距也超过10 mm；
两者数量级不同。

\thermotransportproblem{228}{88}
标准状况下氮气，求：
(a) 每个分子每秒平均碰撞次数；
(b) $1\,\mathrm{cm^3}$ 氮气中每秒
所有分子间的碰撞总次数。
\exsolution
分子数密度
$n=p_0/(k_BT_0)$，
平均速率
$\bar v=\sqrt{8RT_0/(\pi M)}$。
每分子平均碰撞率
\[
\boxed{z_1=\frac{\bar v}{\lambda}
=\sqrt2\pi d^2 n\bar v}.
\]
采用氮气有效硬球直径约
$d\simeq0.34$--$0.37\,\mathrm{nm}$，
得到
$z_1$ 为
$10^{10}\,\mathrm{s^{-1}}$ 量级；
原书采用的碰撞截面取值给
$\boxed{z_1\simeq0.74\times10^{10}\,\mathrm{s^{-1}}}$。
(b) 每次碰撞涉及两个分子，
计算全部碰撞对时须除以2。
体积 $V=1\,\mathrm{cm^3}$ 中分子数为 $nV$：
\[
\boxed{Z=\frac12nVz_1}.
\]
取 $n\simeq2.7\times10^{19}\,\mathrm{cm^{-3}}$：
\[
\boxed{Z\simeq1.0\times10^{29}\,
\mathrm{s^{-1}}}
\quad\text{（每立方厘米）}.
\]
\excheck 碰撞总次数不能简单写成 $nVz_1$，
否则每个碰撞对会被重复计数两次。

\thermotransportproblem{229}{89}
在理想气体的
(a) 等容过程；(b) 等压过程中，
分子平均自由程 $\lambda$
与每分子单位时间碰撞次数 $z$
如何依赖绝对温度 $T$？
\exsolution
硬球模型中
$\lambda=(\sqrt2\pi d^2n)^{-1}$，
平均分子速率 $\bar v\propto\sqrt T$，
碰撞次数
$z=\bar v/\lambda$。
(a) 等容且物质的量不变，
数密度 $n=N/V$ 不变，
所以
\[
\boxed{\lambda=\text{常数},\qquad z\propto T^{1/2}}.
\]
(b) 等压过程中理想气体
$n=p/(k_BT)\propto T^{-1}$，
故
\[
\boxed{\lambda\propto T,\qquad z\propto T^{-1/2}}.
\]
\excheck 不可把分子平均速率误当成 $\propto T$；
麦克斯韦平均速率正比于 $\sqrt T$，
这决定碰撞频率的半整数幂。
